%==============================================================================
\section{Hypothesis --- Giả thuyết trong học máy}
\label{sec:hypothesis}
%==============================================================================

\begin{dinhnghia}[title={Hypothesis (giả thuyết)}]
Trong ML, giả thuyết là \textbf{lời giải thích hoặc giải pháp đề xuất} cho bài toán --- giả định tạm thời có thể kiểm tra bằng dữ liệu.
Trong học có giám sát, giả thuyết chính là \textbf{mô hình} thuật toán học để dự đoán trên dữ liệu mới.
\end{dinhnghia}

\begin{ynghia}[title={Ví dụ đời thường}]
``Giá nhà tỷ lệ thuận với diện tích sàn'' --- một giả thuyết có thể kiểm tra bằng dữ liệu.
\end{ynghia}

\subsection{Giả thuyết trong Machine Learning}

\begin{ynghia}[title={Hàm giả thuyết}]
Giả thuyết thường là hàm ánh xạ đầu vào sang dự đoán:
\[
h(x) = \hat{y}
\]
với $x$ là đầu vào, $\hat{y}$ là giá trị dự đoán.
target ML: tìm giả thuyết \textbf{tổng quát hoá tốt} trên dữ liệu chưa thấy.
\end{ynghia}

\subsection{Hàm giả thuyết $h$}

\begin{dinhnghia}[title={Hypothesis function}]
Hàm giả thuyết nhận đầu vào và trả đầu ra.
Hồi quy tuyến tính đơn biến:
\[
h(x) = w_0 + w_1 x
\]
$w_0$, $w_1$ là tham số (trọng số); $x$ là feature.

Hồi quy tuyến tính đa biến:
\[
h(x) = w_0 + w_1 x_1 + \cdots + w_n x_n
\]
\end{dinhnghia}

\begin{phuongphap}[title={Huấn luyện}]
Quá trình tìm giá trị tối ưu cho tham số sao cho \textbf{hàm mất mát} (cost) nhỏ nhất.
\end{phuongphap}

\subsection{Không gian giả thuyết $\mathcal{H}$}

\begin{dinhnghia}[title={Hypothesis space}]
Tập mọi giả thuyết có thể có gọi là \textbf{không gian giả thuyết}.
Với hồi quy tuyến tính, không gian gồm mọi hàm tuyến tính có thể.
Thuật toán tìm giả thuyết ``khớp'' nhất trong không gian này --- gọi là \textbf{huấn luyện / học}.
\end{dinhnghia}

\subsection{Hai loại giả thuyết (kiểm định thống kê)}

\begin{dinhnghia}[title={Null hypothesis ($H_0$)}]
Giả thuyết \textbf{mặc định}: không có quan hệ giữa feature và biến target.
Trong ML ta thường cố \textbf{bác bỏ} $H_0$ nếu $p$-value $< \alpha$ (mức ý nghĩa).
\end{dinhnghia}

\begin{dinhnghia}[title={Alternative hypothesis ($H_1$)}]
Giả thuyết \textbf{đối lập} $H_0$: có quan hệ có ý nghĩa giữa đầu vào và đầu ra.
Bác bỏ $H_0$ $\Rightarrow$ chấp nhận $H_1$.
\end{dinhnghia}

\subsection{Kiểm định giả thuyết trong ML}

\begin{phuongphap}[title={Các bước}]
\begin{enumerate}
  \item Phát biểu $H_0$ và $H_1$.
  \item Chọn mức ý nghĩa $\alpha$ (thường 0{,}05 hoặc 0{,}01).
  \item Tính thống kê kiểm định ($t$, $z$, \ldots).
  \item Tính $p$-value; nếu $p < \alpha$ thì bác bỏ $H_0$.
  \item Kết luận: quan hệ có ý nghĩa hay không.
\end{enumerate}
\end{phuongphap}

\subsection{Tìm giả thuyết tốt nhất}

\begin{ynghia}[title={Tối ưu hoá}]
Huấn luyện = điều chỉnh tham số để tối thiểu hóa loss (sai khác giữa $\hat{y}$ và $y$).
Kỹ thuật như \textbf{gradient descent} tìm $w$ tối ưu.

Ví dụ hồi quy tuyến tính, cost (MSE):
\[
J(w) = \frac{1}{2n}\sum_{i=1}^{n}\bigl(h(x_i) - y_i\bigr)^2
\]
target: $w$ làm $J(w)$ nhỏ nhất $\Rightarrow$ giả thuyết tốt nhất trong lớp hàm đã chọn.
\end{ynghia}

\subsection{Tính chất giả thuyết tốt}

\begin{tinhchat}[title={Một giả thuyết / mô hình tốt nên}]
\begin{itemize}
  \item \textbf{Generalization}: dự đoán tốt trên dữ liệu mới.
  \item \textbf{Simplicity}: đơn giản, dễ diễn giải.
  \item \textbf{Robustness}: chịu được nhiễu và outlier.
  \item \textbf{Scalability}: xử lý được dữ liệu lớn.
\end{itemize}
\end{tinhchat}

\begin{phuongphap}[title={Các học thuật toán sinh giả thuyết}]
Linear/logistic regression, cây quyết định, SVM, mạng nơ-ron, \ldots

Sau huấn luyện cần \textbf{đánh giá} trên tập validation hoặc cross-validation trước khi triển khai thực tế.
\end{phuongphap}
